\section{Experiments}
\label{sec:experiments}
\subsection{Platform and task}
We evaluate the fixed \texttt{paper1-eval-finite-scan-v1} system in Gazebo
physics-based simulation. P450 supplies aerial RGB-D and MID360 observations;
BUNKER carries an AUBO i5 arm, an AG95 gripper and a wrist-mounted D435 camera.
The shared stack uses PX4 SITL/Prometheus for aerial control, ROS navigation
for the base, and MoveIt-based arm planning. Each task begins with normal aerial
target perception and proceeds through the policy's own observations, base
selection and robot execution. No successful base pose is manually assigned.
The perceived target is a known brick on static horizontal ground. Physical
grasping uses simulated contact and normal controller execution, not an object
teleportation or a substituted success label.

All methods use the same initial observation pose
$(x,y,z,\psi)=(-1.4,0,1.2,0)$ in map, the same RGB-D acceptance gate, aligned
$0.10\,\mathrm{m}$ fields, and the settings of Section~\ref{sec:method} where
applicable. The actual collection revisions are AGENT \texttt{3d13a95},
SIM \texttt{a0ae8e3} and original RM4D \texttt{e9d4312}, with the separately
calibrated Ground-task asset. These identify the executed system rather than
assuming that a later development tip reproduces it. No robot decision rule,
scene seed, sensor gate or physical success predicate changed during collection.

\subsection{Independent scenes and paired protocol}
\label{sec:scenes}
The primary study contains 96 independently sampled scenes. Difficulty is drawn
independently with probability $1/3$ for Easy, Moderate and Hard, giving the
prospectively fixed realization of 38, 28 and 30 scenes. Easy contains no wall,
Moderate one wall, and Hard two walls. The walls are $1\,\mathrm{m}$ high;
the Moderate wall is $1\,\mathrm{m}$ long and Hard walls $0.6\,\mathrm{m}$ long,
all with $0.15\,\mathrm{m}$ thickness. Target position/yaw, initial base pose,
and wall position/yaw are independently perturbed according to the frozen
generation rules. The target is centered near $(2,0)$ and the base near
$(3,-2.5)\,\mathrm{m}$. Target XY perturbations are bounded by $0.15\,\mathrm{m}$
and yaw by $\pi/6$; wall-center perturbations are bounded by
$0.15\,\mathrm{m}$ and wall yaw by $\pi/12$.

The Hard generator introduces opportunities for occlusion alongside unknown
space that need not support manipulation; it does not guarantee that each
realization produces either a blocked useful region or a favorable method
difference. Seeds were disjoint from development data, and the complete
geometry, method order and auxiliary subsets were fixed before method runs.
Admissibility checks concern initial geometry and sensor setup, not candidate
counts, path existence, viewpoint scores or retrieval outcomes. Each scene is
run once per primary method in the predeclared balanced/shuffled order, yielding
192 planned primary trials. The paired scene, not the view or candidate, is the
independent analysis unit. We retain the realized mixture rather than
post-stratifying tiers after observing their outcomes.

\subsection{Comparators and ablations}
\label{sec:baselines}
The sole primary contrast is \emph{Generic NBV} versus \emph{Ours}. Generic uses
Eq.~\eqref{eq:genericgain}; Ours uses Eq.~\eqref{eq:taskgain}. They share initial
conditions, candidate generation, finite-scan opportunity and occlusion, cost,
observation budget, evidence updates, exact support, operational gating,
candidate screening and the full execution pipeline. Their later measurements
may differ because their actions differ. Generic is also the without-task-
weighting ablation; it is not counted a second time.

Two context controls run on six predefined scenes, the first two generated
IDs in each tier. \emph{RM4D-only} uses the initial aerial RGB-D target and
calibrated task-domain map, selects the baseline top candidate, and skips
MID360 confirmation and pre-handoff multi-candidate screening. Downstream
physical checks remain enabled. It is therefore a context control, not a
controlled isolation of relevance weighting or a test of the original
literature map's domain. \emph{Fixed-view} reobserves the current measured hover
location without selecting translated NBVs, retaining the three-window cap and
shared screening stop.

On the first four generated Hard IDs, \emph{without occlusion} removes occlusion
only from the predicted gain, not from collision blocking or execution checks;
\emph{without flight cost} sets the policy flight weight to zero. The scheduled
Ours counterpart is reused for each auxiliary comparison. These 12 context and
8 ablation tasks bring the total to 212 planned tasks, not 212 primary trials.
The small auxiliary subsets provide descriptive context and do not support
powered secondary superiority, equivalence or universal-necessity claims.

\subsection{Outcomes, failures and observation cost}
\label{sec:metrics}
The primary binary endpoint is physical retrieval success. It requires the
perception, arrival, grasp and lift/short-retention checks, including an
independent simulated-object height check. The existing checker requires both
TCP lift and measured object-height increase of at least $0.10\,\mathrm{m}$;
the task also requires its existing one-second wall-clock hold of fresh grasp
confirmation. This interface hold is distinct from the simulation-time efficiency
metrics below. Candidate confirmation, a successful
planning preview, $\Dexec$, TCP motion or a task-process success message alone
does not establish retrieval. We record progress through observation,
confirmation, screening, navigation/stopping, D435 refinement, $\Dexec$, grasp
and lift/retention. First-terminal failure counts assign one cause to an
unsuccessful task; downstream stages that were not reached do not create
additional failures. Model rejection is not interpreted as proof that a
physical collision occurred.

The common observation budget is three accepted five-simulation-second windows,
including the initial window. A further observation may follow translation,
yaw-only motion or reobservation near the same location. We distinguish accepted
windows from actual changes between measured observation poses. A resolved
location change exceeds $0.20\,\mathrm{m}$ translation; wrapped yaw change
exceeding $0.20\,\mathrm{rad}$ is separately identifiable. These are reporting
categories, not flight-path lengths or additional control thresholds.

Robot-efficiency times use simulation time. Active time starts at the first
capture and ends at active stop or terminal failure, including observation,
motion, waiting, computation and shared previews. It excludes initial takeoff,
aerial RGB-D and the initial reachability query. Ground time starts at navigation
and ends at the task's terminal Ground event. Whole-task time includes
initialization and the normal return/landing sequence. Unreached Ground stages
have missing Ground time, not zero cost. Distance traces are incomplete in this
release, so distance saving is not an inferential endpoint or claim.

Normal insufficient sensing, bounded planning rejection and physical execution
failure remain valid outcomes. Only documented infrastructure-invalid attempts
may receive one same-scene/method replacement within the fixed reserve. All
212 planned slots have valid selected binary outcomes and all 96 primary pairs
are complete. Two proven-invalid originals received replacements, giving 214
formal activations; one additional external bare simulator startup was separately
charged to resources. Original attempts are preserved. These activations are not
additional independent samples.

\subsection{Prespecified analysis}
\label{sec:analysis}
Let $a,b,c,d$ denote both-success, Ours-only, Generic-only and both-failure paired
counts. We report the marginal success rates and
\begin{equation}
 \widehat\Delta=(b-c)/96.
 \label{eq:effect}
\end{equation}
The single primary test is the two-sided exact McNemar test, conditioning on
$b+c$ discordances, with significance level $0.05$. The prespecified conservative
95\% interval for $\Delta$ combines two 97.5\% Clopper--Pearson intervals
$[L_b,U_b]$ and $[L_c,U_c]$ for $b/96$ and $c/96$ as
$[L_b-U_c,U_b-L_c]$. Bonferroni joint coverage is at least 95\%; independence of
the two discordant categories is not assumed. The interval concerns the paired
risk difference, not the test's $p$-value. A first-activation sensitivity counts
infrastructure-invalid originals as unsuccessful without dropping denominators.

Tier-specific outcomes and auxiliary comparisons are descriptive, not additional
primary tests. Conditional efficiency is summarized on the 53 scenes where both
primary methods retrieve successfully, with 10,000 whole-scene paired bootstrap
resamples and the fixed analysis seed. Conditioning is explicit: these means
are not unbiased all-task cost effects, and a failed quick task is not efficient
retrieval. The existing analysis also reports all-96-denominator success-resource
curves. No additional experiment, significance test or outcome-driven sample
extension is introduced in this manuscript.
